% Welcome to the unofficial template of the master study Smart Systems Engineering at Hanze University of Applied Sciences. Original template created by Natanael Magno Gomes and The Polytechnic Institute of Bragança (IPB).
% Adapted to Hanze's expectations by dr. Felipe Nascimento Martins and subsequently by Ewout Bergsma.
% All information of the "05_Final_report_format_20_21.docx" document is included in this template.
% Last update: 12/12/2021
% Feel free to contact me: ewoutbergsma [at] hotmail.com

\documentclass[12pt,a4paper,twoside]{ipb}
\errorcontextlines=3
\usepackage{csquotes}
\usepackage{amssymb}
\usepackage[portuguese,USenglish]{babel}
\graphicspath{{./images/}}
\usepackage{listings}
\usepackage{xcolor}
\usepackage{realboxes}
\usepackage{url}
\usepackage[colorlinks=true, urlcolor=black, linkcolor=black, citecolor=black, bookmarks=true, pdfstartview=FitH]{hyperref}

\usepackage[style=ieee,backend=biber]{biblatex}
\addbibresource{references.bib}
\usepackage{lipsum}
\usepackage{tabulary}
\usepackage{graphicx}
\usepackage{subfig}
\usepackage{float}  % Positioning of the figures
\usepackage[binary-units=true]{siunitx}
\usepackage[section]{placeins}
\usepackage{indentfirst}  % Indents all paragraphs
\usepackage{mathtools}  % Allows adding multi line equations
\emergencystretch=1em  % Fix overfull
\newcolumntype{P}[1]{>{\centering\arraybackslash}p{#1}}  % Centre col with width
\usepackage{enumitem}
\setlist{nosep}  % By default the separation is large
\renewcommand{\glsgroupskip}{}


\usepackage{fancyref}  % This allows to use \fref{} command. This is very useful, it allows to easily refer to any \label{}. % See documentation: https://mirror.las.iastate.edu/tex-archive/macros/latex/contrib/fancyref/fancyref.pdf.


\title{Unofficial template Master SSE Thesis}  % Title of your thesis
\author{Studentname Middlename Lastname}  % Student name
\supervisor{Name First Supervisor}
% \cosupervisor{Name co-supervisor (Uncomment this line if there is no co-supervisor)}
\courseyear{2021}
\makeglossaries
\loadglsentries{chapters/acronyms}


\begin{document}
\beforepreface  % Places the cover page

\clearpage
% Navigate to: chapters/abstract.tex
\begin{center}
    ABSTRACT
\vspace{5mm} %5mm vertical space
\end{center}

%%%%%%%%%%%%%%%%%%%%%%%%%%%%%%
Insert your abstract here. Original Hanze format states one can use a maximum of 200 words. Though, it is advised is to discuss with your supervisor(s) if you need more.
%%%%%%%%%%%%%%%%%%%%%%%%%%%%%%

\vspace{5mm} %5mm vertical space
\noindent {\bf Keywords:} Keyword, Keyword, Keyword.  % Replace keywords


\clearpage
% Navigate to: chapters/declaration.tex
\begin{center}
    DECLARATION
    \vspace{5mm} %5mm vertical space
\end{center} 

I hereby certify that this report constitutes my own product, that where the language of others is set forth, quotation marks so indicate, and that appropriate credit is given where I have used the language, ideas, expressions or writings of another.

\vspace{1em}

I declare that the report describes original work that has not previously been presented for the award of any other degree of any institution.

\vspace{2em}

\noindent\hspace{8cm}Signed,
\begin{figure}[H]
    \hspace{8cm}
    \includegraphics[width=6cm]{signature.png}  % Place an image of your signature in location: images/signature.png
    % Note how figures without \caption{} do not appear in the list of figures (same holds for tables)
\end{figure}

\noindent\hspace{8cm}Studentname Middlename Lastname  % Replace with your name


\clearpage
% Navigate to: chapters/acknowledgements.tex
\begin{center}
    ACKNOWLEDGEMENTS
    \vspace{5mm} %5mm vertical space
\end{center} 

Here you have the opportunity to thank anyone. Acknowledgements should be in good taste and should not extend more than one page.


\cleardoublepage  % Makes sure contents start on correct page
\afterpreface  % Places contents, list of tables and list of figures
\bodystart  % This sets stylistic parameters for the following content

% Here starts the actual content of your thesis
\chapter{Rationale}\label{chap:rationale}
The first chapter of the Final Report should explain the rationale. It can be the same chapter as in your Master thesis definition, but you should incorporate any feedback that was given to you by your first graduation supervisor and your company supervisor.

\chapter{Situational \& Theoretical Analysis}\label{chap:situational_theoretical_analysis}
This chapter in your final report can also be the same as in your master thesis definition, but should include all feedback that was given to you.
\chapter{Conceptual Model}\label{chap:conceptual_model}

In this chapter you must explain, using the information from the previous two chapters and critical thinking, the most likely solution to your problem definition (narrowing down the theories and research), resulting in a final theory and/or model. Under a conceptual model it is understood that:  the author first investigates which factors play a role, and then provides argumentation for the identification of the most important ones. This chapter is basically a summary of chapter 2 and should be roughly 1 page long.



\chapter{Research design}\label{chap:research}

The Research Design chapter contains everything that you have done in order to answer your research question. Write and state things in such a way that somebody can repeat your research and arrive at the exact same results. You can use the chapter in your master thesis definition as a start, but generally, research designs end up differently from what you intended to do in the first place. A special remark: it should be written in the PAST TENSE, since you now have already done it.


Note, the following content are some examples of techniques that one can use. Read the comments for explanations.


\section{Methodology}

A list of requirements has been defined to guide the project, a general and detailed design were developed and implemented to support the investigation of the research question.

\subsection{Requirements}

% One can replace \begin{enumerate} and \end{enumerate} with \begin{itemize} and \end{itemize} respectively for bullets instead of numbers.
\vspace{2mm}
\begin{enumerate}
  \item The Cobot safety features are enabled at all times.
  \item The system architecture is ROS based.
  \item The system works in simulation and real environment.
  \item The training is done using RL techniques.
  \item The algorithm is trained using self experience and memory replay.
  \item The algorithm uses RGBD images to estimate the actions.
\end{enumerate}
\vspace{3mm}

\section{General Design}

General design text. See \fref{tab:models}.  % Note how \fref{} allows to refer to a label.

\section{Detailed Design}

Detailed design text. You don't have to follow this structure. Feel free to include, remove or change the titles of the sections.

An example of a table is shown in \fref{tab:models}. Moreover, an example of a figure is shown in \fref{fig:example_figure}. Always use the proper reference to cite tables and figures. Do not mention "table below" or "figure above", for example.

\begin{table}
    \centering
        \begin{tabular}{ p{35mm}  P{3cm}   P{40mm}  P{3cm}}
            \hline
            Model name & Feature space & Est. total size (MB) & Top-5 error \\
            \hline
            AlexNet          &   256  x 6 x  6     &  242.03  &    20.91 \\
            SqueezeNet       &   512  x 13 x 13   &   76.54   &    19.58 \\
            Densenet         &   1024 x 7 x  7    &   325.21  &    7.83  \\
            GoogleNet        &   1024 x 7 x  7    &   85.91   &    10.47 \\
            MobileNet        &   1280 x 7 x  7    &   162.45  &    9.71  \\
            ResNeXt-50-32x4d &   2048 x 7 x  7    &   415.72  &    6.30  \\
            MNASNet          &   1280 x 7 x  7    &   135.77  &    8.4  \\
            \hline
        \end{tabular}
    \caption{This is the caption of the example table \cite{torchvisionmodels}.}
    \label{tab:models}
\end{table}


\begin{figure}[b]  % b for bottom of page
    \centering
    \includegraphics[width=0.8\textwidth]{images/hanzelogo_nl.png}
    \caption{This is the caption of the example figure.}
    \label{fig:example_figure}
\end{figure}

\chapter{Research Results}\label{chap:research_results}

In the chapter discussing your research results you must provide all the results of your research in tables, graphs and figures format. Compare your results with the results found in literature or previous research and include all falsifications that you are aware of. Also comparisons must be made with the hypothesis you stated in the second chapter: are they correct? If so, why? If not, why not?
This chapter should include all the information that came out of your research and it should include a critical interpretation of these results.


\chapter{Conclusions and Recommendations}\label{cap:conclusions_recommendations}

The last chapter contains the conclusions and recommendations. All your conclusions should be a logical consequence of the previous chapters that you have written and should provide answers to the research questions and problem definition posed in chapter \ref{chap:rationale}, \ref{chap:situational_theoretical_analysis} and \ref{chap:conceptual_model} 



\chapter{Additional Chapters}\label{chap:additional_chapters}

If you want, you can add additional chapters to your final report. Keep in mind that the mentioned chapters should at least be present and the main text (excluding title page, table of contents, references and appendices) of the final report should not be more than 15000 words.

Note, if you need more words, please discuss with your supervisor(s) \ref{ISO15066}.

\begin{enumerate}
    \item \gls{AI}
    \item \gls{AI}
    \item \Gls{AI}
    \item \glspl{AI}
    \item \acrshort{AI}
    \item \acrlong{AI}
    \item \acrfull{AI}
\end{enumerate}


\printglossary[type=\acronymtype,title={Definitions and Abbreviations}]

\clearpage
\addcontentsline{toc}{chapter}{References}
\printbibliography[title={References}]

\chapter*{Appendix}\label{appendix}
% Following is required as above uses \chapter*{} (note the star). The start makes the chapter unnumbered, but also removes it from table of content. Former is desired, the latter is not:
\addcontentsline{toc}{chapter}{Appendix}


You are encouraged to put in appendices in your final report. In an appendix you can include things such as large tables or background information. Anything that is useful to know for the reader, but prevents the reader to read your main text in a fluent manner. Each appendix should have a number and a self-explanatory title.  % Note, appendix must be last

\end{document}
